% =========================================================
% Covers, Subcovers, and the Finite Intersection Property
% =========================================================

\subsection{Covers, Subcovers, and the Finite Intersection Property}

% ---------------------------------------------------------
% TOOLKIT
% ---------------------------------------------------------

\begin{exposition}
Compactness is one of the central concepts of analysis, and its definition is
stated in terms of covers: a space is compact if every open cover has a finite
subcover. This is a purely set-theoretic condition --- it makes no reference to
distance or continuity. The work of understanding compactness therefore begins
here, with the vocabulary of covers. Once these definitions are in place, the
definition of compactness will be a single sentence, and the difficulty will
lie entirely in the geometric and analytic content, not the set-theoretic
framework.
\end{exposition}

% ---------------------------------------------------------
% Covers
% ---------------------------------------------------------

\begin{definitionbox}{Definition (Cover)}
\begin{definition}[Cover]\label{def:cover-full}
Let $A$ be a set and $\mathcal{C}$ a collection of sets.
$\mathcal{C}$ is a \emph{cover} of $A$ if every element of $A$ belongs to
at least one member of $\mathcal{C}$:
\[
A \;\subseteq\; \bigcup_{C \in \mathcal{C}} C.
\]
\end{definition}
\end{definitionbox}

\begin{remark*}[Standard quantified statement]
\[
  \mathcal{C}\text{ covers }A
  \iff
  \forall x\in A,\ \exists C\in\mathcal{C},\ x\in C.
\]
\end{remark*}

\begin{remark*}[Predicate reading]
\[
  \operatorname{Covers}(\mathcal{C},A)
  \iff
  \forall x\in A,\ \exists C\in\mathcal{C},\ x\in C.
\]
\end{remark*}

\begin{remark*}[Interpretation]
A cover supplies at least one set containing each element of the covered set.
\end{remark*}

\LeanFormalizes{def:cover-full}{lra-lean}{LRA.VolumeI.Set.Families}{Covers}{pending}
\LeanFormalizes{def:cover-full}{lra-lean}{LRA.VolumeI.Set.Families}{CoversElementwiseIff}{pending}

\begin{dependencies}
\begin{itemize}
  \item \hyperref[def:subset]{Subset}
  \item \hyperref[def:union]{Union}
\end{itemize}
\end{dependencies}

\begin{definitionbox}{Definition (Subcover)}
\begin{definition}[Subcover]\label{def:subcover}
$\mathcal{C}'$ is a \emph{subcover} of $\mathcal{C}$ if $\mathcal{C}'
\subseteq \mathcal{C}$ and $\mathcal{C}'$ is still a cover of $A$.
\end{definition}
\end{definitionbox}

\begin{remark*}[Standard quantified statement]
\[
  \mathcal{C}'\text{ is a subcover of }\mathcal{C}\text{ for }A
  \iff
  \mathcal{C}'\subseteq\mathcal{C}
  \wedge
  \operatorname{Covers}(\mathcal{C}',A).
\]
\end{remark*}

\begin{remark*}[Interpretation]
A subcover is a smaller collection that still covers the same set.
\end{remark*}

\LeanFormalizes{def:subcover}{lra-lean}{LRA.VolumeI.Set.Families}{Subcover}{pending}
\LeanFormalizes{def:subcover}{lra-lean}{LRA.VolumeI.Set.Families}{SubcoverIff}{pending}

\begin{dependencies}
\begin{itemize}
  \item \hyperref[def:cover-full]{Cover}
  \item \hyperref[def:subset]{Subset}
\end{itemize}
\end{dependencies}

\begin{definitionbox}{Definition (Finite Cover)}
\begin{definition}[Finite Cover]\label{def:finite-cover}
A cover is \emph{finite} if it has finitely many members. A \emph{finite
subcover} is a subcover $\mathcal{C}' \subseteq \mathcal{C}$ with
$|\mathcal{C}'| < \infty$.
\end{definition}
\end{definitionbox}

\begin{remark*}[Standard quantified statement]
\[
  \mathcal{C}'\text{ is a finite subcover of }\mathcal{C}\text{ for }A
  \iff
  \mathcal{C}'\subseteq\mathcal{C}
  \wedge
  \operatorname{Covers}(\mathcal{C}',A)
  \wedge
  \operatorname{FiniteSet}(\mathcal{C}').
\]
\end{remark*}

\begin{remark*}[Interpretation]
A finite subcover covers the same set using only finitely many members of the
original cover.
\end{remark*}

\LeanFormalizes{def:finite-cover}{lra-lean}{LRA.VolumeI.Set.Families}{FiniteCollection}{pending}
\LeanFormalizes{def:finite-cover}{lra-lean}{LRA.VolumeI.Set.Families}{FiniteSubcover}{pending}
\LeanFormalizes{def:finite-cover}{lra-lean}{LRA.VolumeI.Set.Families}{FiniteSubcoverIff}{pending}

\begin{remark*}[Exposition]
The condition of being a cover is easy to satisfy: take $\mathcal{C} =
\mathcal{P}(A)$ and it is trivially a cover. The interesting question is
whether large covers can be replaced by smaller ones --- specifically, finite
ones. An infinite cover has infinitely many sets, which is hard to work with.
A finite subcover is just as good for covering purposes, but is finite and
therefore tractable. When a finite subcover always exists, no matter which
cover one starts with, the space is called compact.
\end{remark*}

\begin{dependencies}
\begin{itemize}
  \item \hyperref[def:cover-full]{Cover}
  \item \hyperref[def:finite-infinite]{Finite and Infinite Sets}
  \item \hyperref[def:subset]{Subset}
\end{itemize}
\end{dependencies}

\begin{definitionbox}{Definition (Open Cover)}
\begin{definition}[Open Cover]\label{def:open-cover}
In a topological space $(X, \tau)$, an \emph{open cover} of a subset $A
\subseteq X$ is a cover $\mathcal{C}$ of $A$ all of whose members are open
sets: $\mathcal{C} \subseteq \tau$.
\end{definition}
\end{definitionbox}

\begin{remark*}[Standard quantified statement]
$\mathcal{C}$ is an open cover of $A$ $\;\leftrightarrow\;$ $(\forall C \in \mathcal{C},\; C \in \tau) \;\wedge\; (\forall x \in A,\; \exists C \in \mathcal{C},\; x \in C)$.
\end{remark*}

\begin{remark*}[Predicate reading]
\[
  \operatorname{OpenCover}(\mathcal{C},A,(X,\tau))
  \iff
  \operatorname{Covers}(\mathcal{C},A)
  \wedge
  \forall C\in\mathcal{C},\ C\in\tau.
\]
\end{remark*}

\begin{remark*}[Interpretation]
An open cover is a cover whose covering sets are open in the ambient
topological space.
\end{remark*}

\begin{remark*}[Exposition]
The emphasis on open covers in compactness is because open sets are the
sets one controls in a topological space. An open cover $\{U_\alpha\}$ of
$A$ means: every point of $A$ has at least one $U_\alpha$ containing it. A
finite subcover then gives a list of finitely many open sets whose union
already contains all of $A$. Compactness says this reduction is always
possible, which is a very strong structural property.
\end{remark*}

\begin{dependencies}
\begin{itemize}
  \item \hyperref[def:cover-full]{Cover}
  \item \hyperref[def:fip]{Finite Intersection Property}
  \item \hyperref[def:complement]{Complement}
  \item \hyperref[def:intersection]{Intersection}
  \item \hyperref[def:subset]{Subset}
\end{itemize}
\end{dependencies}

% ---------------------------------------------------------
% Finite Intersection Property
% ---------------------------------------------------------

\begin{definitionbox}{Definition (Finite Intersection Property)}
\begin{definition}[Finite Intersection Property]\label{def:fip}
A collection $\mathcal{F}$ of sets has the \emph{finite intersection
property} (FIP) if every finite subcollection has nonempty intersection:
for all finite $\mathcal{F}' \subseteq \mathcal{F}$,
\[
\bigcap_{F \in \mathcal{F}'} F \;\neq\; \varnothing.
\]
\end{definition}
\end{definitionbox}

\begin{remark*}[Standard quantified statement]
\[
  \mathcal{F}\text{ has the FIP}
  \iff
  \forall \mathcal{F}'\subseteq\mathcal{F},\
  \operatorname{FiniteSet}(\mathcal{F}')
  \to
  \bigcap_{F\in\mathcal{F}'}F\neq\varnothing.
\]
\end{remark*}

\begin{remark*}[Predicate reading]
\[
  \operatorname{FiniteIntersectionProperty}(\mathcal{F})
  \iff
  \forall \mathcal{F}'\subseteq\mathcal{F},\
  \operatorname{FiniteSet}(\mathcal{F}')
  \to
  \bigcap_{F\in\mathcal{F}'}F\neq\varnothing.
\]
\end{remark*}

\begin{remark*}[Interpretation]
The finite intersection property says that no finite subcollection has empty
intersection.
\end{remark*}

\LeanFormalizes{def:fip}{lra-lean}{LRA.VolumeI.Set.Families}{FiniteIntersectionProperty}{pending}

\begin{remark*}[Exposition]
The FIP is a consistency condition: no finite list of sets from $\mathcal{F}$
is mutually exclusive. It does \emph{not} say the entire intersection
$\bigcap_{F \in \mathcal{F}} F$ is nonempty --- only that every \emph{finite}
sub-intersection is. The difference matters when $\mathcal{F}$ is infinite.

Example: In $\mathbb{R}$, the family $\mathcal{F} = \{(0, 1/n) : n \in
\mathbb{N}\}$ has the FIP (any finite subcollection has nonempty intersection,
since $(0, 1/n) \supseteq (0, 1/N)$ for the largest $N$ in the subcollection).
But $\bigcap_{n=1}^\infty (0,1/n) = \varnothing$. Compact spaces are
characterized by the property that the FIP implies nonempty total
intersection.
\end{remark*}

\begin{dependencies}
\begin{itemize}
  \item \hyperref[def:intersection]{Intersection}
  \item \hyperref[def:empty-set]{Empty Set}
  \item \hyperref[def:subset]{Subset}
  \item \hyperref[def:finite-infinite]{Finite and Infinite Sets}
\end{itemize}
\end{dependencies}

\begin{theorembox}{Theorem (Cover Failure and Relative Complements)}
\begin{theorem}[Cover Failure and Relative Complements]\label{thm:cover-failure-relative-complements}
Let $A$ be a set and let $\mathcal{C}$ be a collection of sets. Then
$\mathcal{C}$ fails to cover $A$ if and only if the relative complements
$A\setminus C$ have nonempty total intersection:
\[
A\not\subseteq\bigcup_{C\in\mathcal C}C
\iff
\bigcap_{C\in\mathcal{C}}(A\setminus C)\neq\varnothing.
\]
\hyperref[prf:cover-failure-relative-complements]{Go to proof.}
\end{theorem}
\end{theorembox}

\begin{remark*}[Standard quantified statement]
\[
A\not\subseteq\bigcup_{C\in\mathcal C}C
\iff
\exists x\in A,\ \forall C\in\mathcal C,\ x\notin C
\iff
\bigcap_{C\in\mathcal C}(A\setminus C)\neq\varnothing.
\]
\end{remark*}

\begin{remark*}[Predicate reading]
\[
  A\not\subseteq\bigcup_{C\in\mathcal C}C
  \iff
  \bigcap_{C\in\mathcal C}(A\setminus C)\neq\varnothing.
\]
The predicate-level content is the equivalence between cover failure and
nonempty total intersection of relative complements.
\end{remark*}

\begin{remark*}[Interpretation]
A cover fails exactly when there is a point of $A$ missed by every member of
the collection.  Such a point belongs to every relative complement
$A\setminus C$.
\end{remark*}

\LeanFormalizes{thm:cover-failure-relative-complements}{lra-lean}{LRA.VolumeI.Set.Families}{RelativeComplementCollection}{pending}
\LeanFormalizes{thm:cover-failure-relative-complements}{lra-lean}{LRA.VolumeI.Set.Families}{CoverFailureIffRelativeComplementIntersectionNonempty}{pending}

\begin{dependencies}
\begin{itemize}
  \item \hyperref[def:cover-full]{Cover}
  \item \hyperref[def:set-difference]{Set Difference}
  \item \hyperref[def:indexed-intersection]{Indexed Intersection}
  \item \hyperref[def:empty-set]{Empty Set}
  \item \hyperref[def:subset]{Subset}
\end{itemize}
\end{dependencies}

\begin{theorembox}{Theorem (No Finite Subcover and FIP)}
\begin{theorem}[No Finite Subcover and FIP]\label{thm:no-finite-subcover-fip}
Let $A$ be a set and let $\mathcal{C}$ be a collection of sets. Then
$\mathcal{C}$ has no finite subcover of $A$ if and only if the collection of
relative complements $\{A\setminus C:C\in\mathcal{C}\}$ has the finite
intersection property:
\[
\bigl(\forall\mathcal{C}_0\subseteq\mathcal{C},\
\mathcal C_0\text{ finite}\to
A\not\subseteq\bigcup_{C\in\mathcal C_0}C\bigr)
\iff
\{A\setminus C:C\in\mathcal{C}\}\text{ has the finite intersection property}.
\]
\hyperref[prf:no-finite-subcover-fip]{Go to proof.}
\end{theorem}
\end{theorembox}

\begin{remark*}[Standard quantified statement]
\[
\begin{aligned}
&\forall\mathcal C_0\subseteq\mathcal C,\ 
\mathcal C_0\text{ finite}\to
A\not\subseteq\bigcup_{C\in\mathcal C_0}C\\
&\qquad\iff
\forall\mathcal C_0\subseteq\mathcal C,\ 
\mathcal C_0\text{ finite}\to
\bigcap_{C\in\mathcal C_0}(A\setminus C)\neq\varnothing.
\end{aligned}
\]
\end{remark*}

\begin{remark*}[Predicate reading]
\[
  \forall\mathcal C_0\subseteq\mathcal C,\ 
  \mathcal C_0\text{ finite}\to
  A\not\subseteq\bigcup_{C\in\mathcal C_0}C
\]
is equivalent to the finite-intersection-property condition for
\(\{A\setminus C:C\in\mathcal C\}\).
\end{remark*}

\begin{remark*}[Interpretation]
The previous theorem turns one failed cover into one nonempty intersection of
relative complements.  This theorem applies that conversion uniformly to every
finite subcollection.  That is exactly the finite-intersection-property
translation used in compactness.
\end{remark*}

\LeanFormalizes{thm:no-finite-subcover-fip}{lra-lean}{LRA.VolumeI.Set.Families}{HasNoFiniteSubcover}{pending}
\LeanFormalizes{thm:no-finite-subcover-fip}{lra-lean}{LRA.VolumeI.Set.Families}{NoFiniteSubcoverIffRelativeComplementFIP}{pending}

\begin{dependencies}
\begin{itemize}
  \item \hyperref[def:cover-full]{Cover}
  \item \hyperref[def:subcover]{Subcover}
  \item \hyperref[def:finite-cover]{Finite Cover}
  \item \hyperref[def:fip]{Finite Intersection Property}
  \item \hyperref[thm:cover-failure-relative-complements]{Cover Failure and Relative Complements}
\end{itemize}
\end{dependencies}

% ---------------------------------------------------------
% FIP--Cover duality
% ---------------------------------------------------------

\begin{propositionbox}{Proposition (FIP--Cover Duality)}
\begin{proposition}[FIP--Cover Duality]\label{prop:fip-duality}
Let $X$ be a set, $A \subseteq X$, and $\{U_\alpha\}_{\alpha \in I}$ a
family of subsets of $X$. Let $F_\alpha = U_\alpha^c = X \setminus U_\alpha$.
Then:
\[
\{U_\alpha\} \text{ covers } A
\;\;\Longleftrightarrow\;\;
\{F_\alpha \cap A\}_{\alpha \in I} \text{ does \emph{not} have the FIP}.
\]
More precisely: the family $\{U_\alpha\}$ does \emph{not} cover $A$ if and
only if the family of closed complements $\{F_\alpha\}$ has the FIP relative
to $A$.
\hyperref[prf:fip-duality]{Go to proof.}
\end{proposition}
\end{propositionbox}

\begin{remark*}[Standard quantified statement]
$A \subseteq \bigcup_{\alpha \in I} U_\alpha \;\leftrightarrow\; \neg(\forall\text{ finite }I_0\subseteq I,\; \bigcap_{\alpha\in I_0} F_\alpha \cap A \neq \varnothing)$. Reading pointwise: $\forall x\in A,\;\exists\alpha\in I,\; x\in U_\alpha \;\leftrightarrow\; \{F_\alpha\cap A\}$ fails the FIP.
\end{remark*}

\begin{remark*}[Predicate reading]
\[
  \operatorname{Covers}(\{U_\alpha\}_{\alpha\in I},A)
  \iff
  \neg\operatorname{FiniteIntersectionProperty}(\{F_\alpha\cap A\}_{\alpha\in I}).
\]
\end{remark*}

\LeanFormalizes{prop:fip-duality}{lra-lean}{LRA.VolumeI.Set.Families}{HasFiniteSubcover}{pending}
\LeanFormalizes{prop:fip-duality}{lra-lean}{LRA.VolumeI.Set.Families}{FiniteSubcoverIffRelativeComplementNotFIP}{pending}

\begin{remark*}[Interpretation]
Covering by a family is equivalent to failure of the finite intersection
property for the relative complements.
\end{remark*}

\begin{remark*}[Exposition]
The duality is a direct consequence of De Morgan's law for indexed families
(Theorem~\ref{thm:indexed-de-morgan}). The family $\{U_\alpha\}$ fails to cover
$A$ means:
\[
A \not\subseteq \bigcup_\alpha U_\alpha
\;\;\Longleftrightarrow\;\;
A \cap \Bigl(\bigcup_\alpha U_\alpha\Bigr)^c \neq \varnothing
\;\;\Longleftrightarrow\;\;
A \cap \bigcap_\alpha U_\alpha^c \neq \varnothing
\;\;\Longleftrightarrow\;\;
A \cap \bigcap_\alpha F_\alpha \neq \varnothing.
\]
Applying this to every finite subcollection gives the equivalence with the FIP.
\end{remark*}

\begin{remark*}[Exposition]
Compactness has two equivalent formulations:
\begin{enumerate}[label=(\roman*)]
\item Every open cover has a finite subcover.
\item Every family of closed sets with the FIP has nonempty intersection.
\end{enumerate}
These formulations are equivalent by FIP--Cover Duality: take complements and
apply De Morgan. Formulation (i) is the standard definition; formulation (ii)
is often easier to apply when working with closed sets (e.g., nested compact
sets). Both are purely set-theoretic conditions, and the equivalence is
entirely a matter of set algebra. When compactness is encountered in topology
and metric spaces, the only new content is the geometric meaning --- the
set-theoretic equivalence is established here.
\end{remark*}

\begin{remark*}[Exposition]
In $\mathbb{R}$, the family $\{[n, \infty) : n \in \mathbb{N}\}$ has the FIP:
any finite subcollection $\{[n_1,\infty), \dots, [n_k,\infty)\}$ has
intersection $[\max(n_i), \infty) \neq \varnothing$. But
$\bigcap_{n=1}^\infty [n, \infty) = \varnothing$. This shows $\mathbb{R}$ is
not compact (the family of closed sets has the FIP but empty total
intersection). A compact space would not allow this failure.
\end{remark*}

\begin{dependencies}
\begin{itemize}
  \item \hyperref[def:cover-full]{Cover}
  \item \hyperref[def:fip]{Finite Intersection Property}
  \item \hyperref[def:complement]{Complement}
  \item \hyperref[def:intersection]{Intersection}
  \item \hyperref[def:subset]{Subset}
\end{itemize}
\end{dependencies}
